% Source: Weinberg GR, physical PDF p. 51
%         (printed p. 25).
% Coverage: Chapter 2 epigraph and introductory paragraphs.
% Figures/tables/footnotes: epigraph; no figures, tables, or footnotes.
% Status: source-reviewed and compile-clean.
% Uncertainties: none.

\begin{flushright}
\begin{minipage}{0.34\textwidth}
\small
``There are really four dimensions, three which we call the three planes of
Space, and a fourth, Time. There is, however, a tendency to draw an unreal
distinction between the former three dimensions and the latter, because it
happens that our consciousness moves intermittently in one direction along the
latter from the beginning to the end of our lives.''

```That,' said a very young man, making spasmodic efforts to relight his cigar
over the lamp; `that \ldots{} very clear indeed.'''

\emph{H. G. Wells, The Time Machine}
\end{minipage}
\end{flushright}

We now review Einstein's Special Theory of Relativity. This chapter, while
self-contained, is only a brief summary, and aims primarily at establishing our
notation and collecting some formulas that will be useful later. The reader who
needs a more extensive introduction to special relativity is advised to turn to
one of the books listed at the end of this chapter, and then return. The reader
who feels completely at home with the subject may find it desirable to move on
immediately to Chapter~3.
